%% Included in decks: slidemage-manual, slidemage-manual-copy
%%

% SlideMage manual, chapter "The big picture": the three parts and the
% workflow, as a described diagram (\tikzillustration; also the example on
% slide 1pn, "Describing diagrams").
\begin{standardframe}{Framework, project and app}
	\small
	SlideMage keeps every slide in a small LaTeX file of its own and builds decks
	from these slides. Three parts work together: the \emph{framework} decides how
	a slide looks, your \emph{project} holds the content, and the \emph{app} does
	the bookkeeping and runs LaTeX for you.
	\par\medskip
	\begin{tikzillustration}[
			part/.style={draw=c_primary_dark, line width=0.6pt, fill=white, rounded corners=2pt,
				text width=108pt, minimum height=62pt, align=left, inner sep=5pt, anchor=north west, font=\footnotesize,
				execute at begin node={\hyphenpenalty=10000\exhyphenpenalty=10000}},
			wstep/.style={fill=c_tertiary_gray, rounded corners=2pt, inner sep=3pt, font=\footnotesize},
		]{1ym-tikz}
		\node[part] (fw) at (0,0) {\textbf{LaTeX framework}\\[2pt]
			\SMCode{latex/}: the Gotham theme, the base preamble, the templates};
		\node[part] (pj) at (150pt,0) {\textbf{Your project}\\[2pt]
			any folder with a \SMCode{slide\_ids.csv}: slides, decks, images};
		\node[part] (app) at (300pt,0) {\textbf{SlideMage}\\[2pt]
			the app: compiles decks, manages IDs, versions, images, transcripts};
		\draw[gotham arrow, c_primary_highlight, thick] (fw.east) -- (pj.west);
		\draw[gotham arrow, c_primary_highlight, thick] (app.west) -- (pj.east);
		% the workflow
		\node[wstep, anchor=west] (s1) at (0,-84pt) {write slides};
		\node[wstep, right=12pt of s1] (s2) {assemble decks};
		\node[wstep, right=12pt of s2] (s3) {compile the deck versions};
		\node[wstep, right=12pt of s3] (s4) {hand out PDFs, transcripts};
		\foreach \a/\b in {s1/s2, s2/s3, s3/s4} \draw[gotham arrow, thick] (\a.east) -- (\b.west);
	\end{tikzillustration}
	\par\medskip
	Compiling means LuaLaTeX through \SMCode{latexmk}, which SlideMage starts for
	you; a bare \SMCode{lualatex} would skip the repeated runs and the bibliography.
\end{standardframe}
